\chapter{Conclusioni e sviluppi futuri}
\label{chap:conclusioni}

Questo capitolo sintetizza il lavoro svolto, riassumendo il problema di partenza, la metodologia adottata e i risultati ottenuti. Vengono inoltre discusse le limitazioni del lavoro di simulazione sviluppato e delineate le possibili direzioni per le ricerche future.

Il settore navale sta attraversando una fase di transizione verso l'impiego di navi a guida autonoma (\ac{MASS}). Tali imbarcazioni necessitano di massicci flussi di dati per mantenere un'adeguata situational awareness, ma le attuali reti di comunicazione \ac{VHF} sono insufficienti a supportare tali carichi~\cite{11217861}. Come punto di partenza, questo lavoro ha analizzato l'algoritmo \emph{Collective Summarisation Clusters} (\ac{CSC}), sviluppato come ``Fase 1" della roadmap CoMPass, il quale utilizza la programmazione aggregata per raggruppare le navi in cluster e trasmettere dati sintetizzati tramite 5G~\cite{11217861, roadmap}. Tuttavia, basandosi esclusivamente su longitudine e latitudine delle navi, questo algoritmo ha mostrato una fragilità topologica in scenari caratterizzati da alta mobilità. 

Per superare questo limite, la presente tesi ha implementato una nuova metrica consapevole della navigazione. Attraverso l'estrazione e la decodifica dei valori \ac{SOG} (velocità) e \ac{COG} (rotta) dai messaggi \ac{AIS}, è stata formulata una funzione di penalità. Questa funzione ha permesso di modificare la scelta dei relay e l'elezione dei leader, scoraggiando la formazione di collegamenti tra imbarcazioni con traiettorie e velocità divergenti. Oltre alle modifiche apportate all'algoritmo di clustering, un contributo rilevante di questo lavoro è consistito nell'estensione del simulatore Alchemist per supportare l'utilizzo completo dei dati \ac{AIS}. In particolare, è stata progettata e implementata una pipeline che consente di decodificare i messaggi \ac{AIS}, estrarre i parametri di navigazione \ac{SOG} e \ac{COG}, mantenerli in memoria durante l'esecuzione della simulazione ed esporli ai programmi Collektive. Questa estensione amplia le capacità del simulatore, costituendo una base riutilizzabile per futuri studi che richiedano l'impiego di ulteriori dati provenienti dai messaggi \ac{AIS}.

I risultati emersi dalla simulazione hanno dimostrato che l'integrazione del criterio di affinità navigazionale influisce sulla topologia della rete. Dal punto di vista prestazionale, il confronto con l'esperimento originario ha evidenziato che l'aggiunta delle variabili cinematiche apporta un miglioramento qualitativo alla stabilità della rete, garantendo al contempo un data rate nettamente superiore rispetto ai metodi tradizionali. Nel confronto tra \ac{M-CSC} e \ac{CSC}, si registra una lieve flessione del data rate complessivo (di circa 100 Kbps) e un andamento temporale più irregolare. Questi dati dimostrano che il minimo costo prestazionale è giustificato dalla creazione di cluster cinematicamente più coerenti e logici. Inoltre, in assenza di infrastruttura 5G ($p_{5G} = 0$), si è registrato un incremento dei singleton clusters (cluster composti da una singola nave) accompagnato da una diminuzione del fattore di riduzione medio. Questi esiti indicano che la sola introduzione delle informazioni di rotta e velocità non è sufficiente, da sola, a innalzare sensibilmente l'efficienza della rete, suggerendo la necessità di strategie di integrazione più complesse.

\section{Limitazioni dello studio}
\label{sec:limitazioni_studio}

Nonostante l'introduzione di dati cinematici rappresenti un passo in avanti rispetto alla sola valutazione della posizione, questo studio è soggetto a diverse limitazioni. In primo luogo, l'esperimento si basa su una quantità di dati modesta, circoscritta alle traiettorie di una finestra temporale di sole sei ore all'interno del fiordo di Kiel, un campione che potrebbe non essere statisticamente rappresentativo del traffico navale su larga scala o in mare aperto~\cite{11217861}. In secondo luogo, l'ambiente di simulazione risulta idealizzato. Il modello non considera la fisica reale che governa il moto delle imbarcazioni (come derive o inerzie), né l'impatto del mutare delle condizioni meteo marine (forza del mare, direzione delle onde o intensità del vento). Dal punto di vista delle telecomunicazioni, il modello assume un canale sostanzialmente affidabile basato sulla distanza, senza simulare fenomeni critici come la perdita di messaggi (\emph{packet loss}), le collisioni sui canali, o il degrado del segnale radio dovuto a ostacoli fisici e avversità atmosferiche. Infine, lo studio non tiene presente anomalie di rete, casi limite (come il guasto improvviso del nodo leader del cluster), né prende in considerazione i problemi aperti legati alla sicurezza visti nella \cref{sec:limiti}~\cite{11217861, 10.1145/2664243.2664257}. Un'ulteriore limitazione riguarda l'impostazione statistica della simulazione. L'esperimento attuale è stato valutato utilizzando un campione di 10 seed rispetto ai 100 impiegati nella simulazione originale. Questo campione ridotto abbassa la stabilità statistica dei risultati, amplificando la deviazione standard e introducendo rumore (evidente nelle oscillazioni dei grafici temporali). Infine, i parametri introdotti per il calcolo della penalità di mobilità (come i pesi attribuiti a \ac{SOG} e \ac{COG}, o la soglia massima di tolleranza di 25 nodi) nel \cref{lst:constants} sono stati determinati euristicamente. A questo si aggiunge che l'algoritmo si limita ad applicare penalizzazioni basate esclusivamente sui valori di velocità e direzione, senza estrapolare informazioni sull'evoluzione futura del sistema.

Alla luce dei risultati e delle limitazioni discusse, i futuri sviluppi di questa ricerca dovranno puntare a simulazioni più ampie e modelli più aderenti alla realtà marittima. Dal punto di vista della valutazione, sarà necessario testare l'algoritmo su dataset \ac{AIS} estesi a più giorni e su aree geografiche differenti. Inoltre, dovrà essere utilizzata la nuova funzionalità di Alchemist discussa nella \cref{sec:integrazione_dati_AIS_alchemist}. Dal punto di vista dell'accuratezza, i simulatori dovranno integrare modelli di propagazione del segnale più complessi, che tengano conto del degrado dovuto al meteo e che introducano la perdita dei messaggi. Per superare i limiti dell'approccio odierno, \ac{SOG} e \ac{COG} potrebbero essere utilizzati per calcolare e prevedere la traiettoria futura delle imbarcazioni. Questo permetterebbe di dotare il sistema di un comportamento di adattamento anticipativo (anticipative adaptation)~\cite{6394122}. Reagendo in modo proattivo alle informazioni locali sugli eventi futuri (come il progressivo allontanamento di un nodo), l'infrastruttura potrebbe calcolare percorsi alternativi in modo emergente (emergent way) e riorganizzarsi prima che i collegamenti si interrompano, sfruttando i concetti dei gradienti anticipativi (\emph{Anticipative Gradients})~\cite{6394122}. Infine, come indicato dalle fasi successive della roadmap CoMPass~\cite{roadmap}, le ricerche future dovrebbero concentrarsi sull'implementazione di algoritmi di machine learning per lo sviluppo di logiche di data summarisation.
